\startchapter[title={Modelos Lineales}, reference=Capitulo_3]

En este capítulo se presentan los dos modelos lineales desarrollados en esta
tesina. Ambos parten del modelo de referencia de Coruzzolo et al.
\cite[coruzzolo2023orderpicking], descrito en la
Sección~\in[sec:trabajos_relacionados]. Primero se define la notación
común a ambos (Sección~\in[sec:notacion]). El modelo monolítico
(Sección~\in[sec:modelo_monolitico]) resuelve ruta, asignación y
secuenciación en un único modelo. El modelo Split
(Sección~\in[sec:modelo_split]) resuelve el mismo problema en dos etapas
sucesivas, con modelos más pequeños, y alcanza el mismo valor óptimo. La
Sección~\in[sec:espera_voluntaria] describe una propiedad común a ambas
formulaciones, la Sección~\in[sec:diferencias_modelos] resume las diferencias
entre ellas y, por último, la Sección~\in[sec:caracteristicas_formulacion] reúne
las características particulares que comparten.

Respecto de ese modelo, la formulación agrega la velocidad heterogénea entre
operarios, la restricción de cadena de frío y un tiempo de servicio en cada
ubicación que atraviesa la ruta. Lo que simplifica es la formación de los
lotes, que aquí son un dato de entrada, de modo que la capacidad del carro
deja de ser una restricción del modelo. Además, la secuencia de cada operario
se expresa con variables de precedencia entre pares de lotes, en lugar de
posiciones numeradas, y el flujo que elimina los subtours se acota con la
cantidad de ubicaciones requeridas del lote, en lugar de la capacidad del
carro (Sección~\in[sec:modelo_monolitico]). La forma de la función objetivo se
mantiene, de modo que, si todos los operarios tienen la misma velocidad, el
límite de la cadena de frío es suficientemente grande y no hay tiempos de
servicio, los dos modelos se reducen al modelo de referencia con los lotes
fijos.

Los dos modelos resuelven el OPP tal como se definió en la
Sección~\in[sec:problema], bajo los siguientes supuestos:

\startitemize
	\item {\bf Lotes dados.} Los lotes se reciben ya conformados y cada orden
	pertenece a un único lote. Como el loteo no se decide, la capacidad del
	carro no forma parte del modelo, y se verifica antes de resolver
	(Sección~\in[sec:split_resolucion]).
	\item {\bf Disponibilidad de los pickers.} Todos los {\em pickers} están
	disponibles desde el instante 0 y ejecutan un lote por vez.
	\item {\bf Efecto de la velocidad.} La velocidad de un {\em picker} escala el
	tiempo de ruta, es decir, el traslado y el servicio en las ubicaciones,
	pero no la preparación del lote ni la recolección de las unidades, que
	dependen de la tarea y no del desplazamiento. Como escala todo el tiempo de
	ruta por igual, la mejor ruta de un lote es la misma para cualquier
	{\em picker}.
	\item {\bf Sin congestión.} Los {\em pickers} no interfieren entre sí al
	recorrer el almacén, de modo que la ruta de un lote tampoco depende del
	momento en que se ejecuta ni de los demás lotes.
	\item {\bf Cadena de frío por lote.} El límite de tiempo fuera de la cámara
	de frío se aplica a la ejecución de cada lote, incluida su parte fija, y
	no al turno del {\em picker} ni a la espera hasta el despacho.
	\item {\bf Extremos de la ruta.} El depósito de partida y el punto de
	entrega son ubicaciones distintas, como en el modelo de referencia.
	\item {\bf Camino simple.} La ruta de un lote no pasa dos veces por la misma
	ubicación, como en el modelo de referencia. Es una simplificación de la
	formulación y no una regla de la operación, y sus consecuencias se
	discuten en la Sección~\in[sec:caracteristicas_formulacion].
\stopitemize

%----------------------------------------------------------------------------------------
\startsection[title={Notación común}, reference=sec:notacion]

Los dos modelos usan los mismos conjuntos y parámetros, que se definen una
sola vez en esta sección.

\startsubsection[title={Conjuntos}]

\startformula
	\startalign[n=2, align={right,left}]
		\NC \mathcal{N} \NC \quad \text{órdenes de clientes} \NR
		\NC \mathcal{B} \NC \quad \text{lotes; cada orden pertenece a un único lote} \NR
		\NC \mathcal{P} \NC \quad \text{{\em pickers}} \NR
		\NC \mathcal{V} \NC \quad \text{ubicaciones del almacén, incluidos el depósito $v_{\text{start}}$ y el nodo final $v_{\text{end}}$} \NR
		\NC \mathcal{E} \NC \quad \text{arcos del grafo de ubicaciones: pares $(i,j)$ tales que se puede ir de $i$ a $j$} \NR
	\stopalign
\stopformula

Para escribir las restricciones de forma compacta se usan además los
siguientes subconjuntos:
\startformula
	\startalign[n=3, align={right,left,left}]
		\NC \mathcal{V}^{\text{int}} \EQ \mathcal{V} \setminus \{v_{\text{start}}, v_{\text{end}}\} \NC \qquad \text{(ubicaciones intermedias)} \NR
		\NC \delta^{+}(i) \EQ \{(i,j) \in \mathcal{E}\} \NC \qquad \text{(arcos que salen de } i\text{)} \NR
		\NC \delta^{-}(i) \EQ \{(j,i) \in \mathcal{E}\} \NC \qquad \text{(arcos que entran a } i\text{)} \NR
		\NC \mathcal{N}_b \EQ \{ n \in \mathcal{N} : B_{nb} = 1 \} \NC \qquad \text{(órdenes del lote } b\text{)} \NR
	\stopalign
\stopformula

\stopsubsection

\startsubsection[title={Parámetros}]

{\bf Órdenes, lotes y almacén.}
\startformula
	\startalign[n=2, align={right,left}]
		\NC A_{ni} \in \{0,1\} \NC \quad \text{1 si la orden $n$ tiene productos en la ubicación $i$} \NR
		\NC B_{nb} \in \{0,1\} \NC \quad \text{1 si la orden $n$ pertenece al lote $b$} \NR
		\NC C_n \NC \quad \text{capacidad (cantidad de unidades a recoger) de la orden $n$} \NR
		\NC T^{\text{due}}_n \NC \quad \text{fecha de entrega ({\em due date}) de la orden $n$} \NR
		\NC D_{ij} \NC \quad \text{distancia del arco $(i,j) \in \mathcal{E}$} \NR
		\NC S_i \NC \quad \text{tiempo de servicio en la ubicación $i$} \NR
	\stopalign
\stopformula

{\bf Tiempos y {\em pickers}.}
\startformula
	\startalign[n=2, align={right,left}]
		\NC T^{\text{travel}} \NC \quad \text{tiempo por unidad de distancia, a velocidad de referencia} \NR
		\NC T^{\text{setup}} \NC \quad \text{tiempo de preparación de cada lote} \NR
		\NC T^{\text{pick}} \NC \quad \text{tiempo de recolección por unidad de capacidad} \NR
		\NC V_p \NC \quad \text{factor de velocidad del {\em picker} $p$ (1,0 = referencia; 2,0 = camina el doble de rápido)} \NR
		\NC T^{\text{cold}} \NC \quad \text{tiempo máximo que un lote puede permanecer fuera de la cámara de frío} \NR
		\NC \varepsilon \NC \quad \text{tolerancia numérica de la restricción de cadena de frío ($10^{-9}$)} \NR
	\stopalign
\stopformula

{\bf Pesos de la función objetivo.}
\startformula
	\startalign[n=2, align={right,left}]
		\NC W^{c} \NC \quad \text{peso de los tiempos de finalización de los lotes} \NR
		\NC W^{\text{pen}} \NC \quad \text{peso del total de penalizaciones} \NR
		\NC W^{e},\; W^{t} \NC \quad \text{pesos del adelanto y de la demora dentro de las penalizaciones} \NR
	\stopalign
\stopformula

{\bf Constantes M grande.}
\startformula
	\startalign[n=2, align={right,left}]
		\NC M \NC \quad \text{constante M grande general, dato de la instancia} \NR
		\NC M^{\text{seq}} \NC \quad \text{constante M grande de las restricciones de secuencia, calculada por instancia} \NR
	\stopalign
\stopformula

\stopsubsection

\startsubsection[title={Parámetros derivados}]

Antes de resolver, a partir de los datos se calculan algunos parámetros
auxiliares. Para cada lote $b$, $a_{bi}$ indica qué ubicaciones debe
visitar: vale 1 si alguna de sus órdenes tiene productos en la ubicación $i$.
A partir de él, $L_b$ cuenta esas ubicaciones. Por último, $F_b$ es el tiempo
fijo del lote, es decir, la preparación más la recolección de todas sus
unidades:
\startformula
	\startalign[n=3, align={right,left,left}]
		\NC a_{bi} \EQ
			\startcases
				\NC 1 \NC \text{si } \sum_{n \in \mathcal{N}_b} A_{ni} \geq 1, \NR
				\NC 0 \NC \text{en caso contrario,} \NR
			\stopcases
		\NC \qquad \text{(1 si alguna orden del lote } b \text{ tiene productos en } i\text{)} \NR
		\NC L_b \EQ \sum_{i \in \mathcal{V}^{\text{int}}} a_{bi} \NC \qquad \text{(cantidad de ubicaciones que el lote } b \text{ debe visitar)} \NR
		\NC F_b \EQ T^{\text{setup}} + T^{\text{pick}} \sum_{n \in \mathcal{N}_b} C_n \NC \qquad \text{(tiempo fijo del lote)} \NR
	\stopalign
\stopformula

{\bf Tiempo de ejecución.} El tiempo fijo $F_b$ no depende de la ruta ni del
{\em picker}. El tiempo que tarda el {\em picker} $p$ en ejecutar el lote $b$
es ese tiempo fijo más el tiempo de ruta ajustado por la velocidad del
{\em picker}:
\startplaceformula[reference=eq:texec]
\startformula
	T^{\text{exec}}_{pb} = F_b + \frac{r_b}{V_p},
\stopformula
\stopplaceformula
donde $r_b$ es el tiempo de ruta del lote (traslado más servicio, a
velocidad de referencia). En el modelo monolítico $r_b$ es una variable de
decisión, así que $T^{\text{exec}}_{pb}$ no es un dato sino una expresión que
depende de la ruta elegida. En el modelo Split, en cambio, $r_b$ proviene de
la primera etapa, y $T^{\text{exec}}_{pb}$ se conoce antes de resolver la
segunda.

{\bf Constante de las restricciones de secuencia.} Las restricciones de
secuencia (Sección~\in[sec:mono_asignacion]) usan una constante M grande
propia, $M^{\text{seq}}$, que se calcula para cada instancia como una cota
del tiempo total de trabajo, es decir, de lo que tardaría el {\em picker} más
lento en ejecutar todos los lotes seguidos, cada uno con una ruta de peor caso
$R^{\max}$. Como una ruta simple recorre a lo sumo $|\mathcal{V}|$ arcos y
pasa por a lo sumo $|\mathcal{V}|$ ubicaciones, $R^{\max}$ acota el tiempo de
cualquier ruta. El resultado se redondea a un entero y se le impone un mínimo
de 500, que solo actúa en instancias muy pequeñas y no compromete la validez
de la cota, porque una constante mayor sigue siendo válida:
\startformula
	\startalign[n=2, align={right,left}]
		\NC M^{\text{seq}} \EQ \max\left(500,\; \sum_{b \in \mathcal{B}} \left( F_b + \frac{R^{\max}}{\min_{p \in \mathcal{P}} V_p} \right) + 1 \right), \NR
		\NC R^{\max} \EQ T^{\text{travel}} \, |\mathcal{V}| \max_{(i,j) \in \mathcal{E}} D_{ij} \;+\; |\mathcal{V}| \max_{i \in \mathcal{V}} S_i. \NR
	\stopalign
\stopformula
Como se vio en la Sección~\in[sec:marco_mip], cuanto menor es una constante
M grande, más ajustada resulta la relajación lineal. Por eso $M^{\text{seq}}$
se calcula a partir de los datos de cada instancia, en lugar de fijarse en un
valor arbitrario. La constante general $M$, en cambio, es un dato de la
instancia.

\stopsubsection
\stopsection

%----------------------------------------------------------------------------------------
\startsection[title={Modelo Monolítico}, reference=sec:modelo_monolitico]

El modelo monolítico es un único \ac{mip} que decide a la vez la ruta de cada
lote, a qué {\em picker} se asigna y en qué orden se ejecuta.

\startsubsection[title={Variables de decisión}]

{\bf Ruteo.}
\startformula
	\startalign[n=2, align={right,left}]
		\NC z_{ijb} \in \{0,1\} \NC \quad \text{1 si la ruta del lote $b$ usa el arco $(i,j)$} \NR
		\NC q_{ijb} \geq 0 \NC \quad \text{ubicaciones requeridas del lote $b$ que quedan por visitar al recorrer el arco $(i,j)$} \NR
		\NC r_b \geq 0 \NC \quad \text{tiempo de ruta del lote $b$ (traslado más servicio)} \NR
	\stopalign
\stopformula

{\bf Asignación y secuencia.}
\startformula
	\startalign[n=2, align={right,left}]
		\NC x_{pb} \in \{0,1\} \NC \quad \text{1 si el lote $b$ se asigna al {\em picker} $p$} \NR
		\NC y_{bk} \in \{0,1\} \NC \quad \text{para $b < k$: 1 si $b$ se ejecuta antes que $k$ (cuando comparten {\em picker})} \NR
	\stopalign
\stopformula

{\bf Tiempos.}
\startformula
	\startalign[n=2, align={right,left}]
		\NC u_{pb} \in [0,\, T^{\text{cold}} + \varepsilon] \NC \quad \text{tiempo de ejecución del lote $b$ si lo ejecuta el {\em picker} $p$} \NR
		\NC c_{pb} \geq 0 \NC \quad \text{instante en que termina el lote $b$ en la agenda del {\em picker} $p$} \NR
		\NC c_b \geq 0 \NC \quad \text{instante en que termina el lote $b$} \NR
		\NC ta_n \geq 0,\; ea_n \geq 0 \NC \quad \text{demora y adelanto de la orden $n$ respecto de $T^{\text{due}}_n$} \NR
	\stopalign
\stopformula

La cota superior de $u_{pb}$ es la que impone la cadena de frío, como se
explica en la Sección~\in[sec:mono_asignacion].

\stopsubsection

\startsubsection[title={Función objetivo}]

Se minimiza la suma ponderada de los tiempos de finalización de los lotes y
de las penalizaciones por adelanto y demora de las órdenes:
\startplaceformula[reference=eq:mono_objetivo]
\startformula
	\min \quad W^{c} \sum_{b \in \mathcal{B}} c_b \;+\; W^{\text{pen}} \left( W^{e} \sum_{n \in \mathcal{N}} ea_n + W^{t} \sum_{n \in \mathcal{N}} ta_n \right)
\stopformula
\stopplaceformula

El tiempo de ruta no aparece explícitamente en el objetivo, sino que influye
a través de los tiempos de finalización y de la cadena de frío.

\stopsubsection

\startsubsection[title={Restricciones de ruteo}]

La ruta de cada lote $b$ es un camino simple desde $v_{\text{start}}$ hasta
$v_{\text{end}}$ que pasa por todas sus ubicaciones requeridas
($a_{bi} = 1$). Puede atravesar otras ubicaciones, pero ninguna más de una
vez.

{\bf Inicio y fin de la ruta.}
La ruta sale una vez del depósito (\in[eq:mono_salida]) y llega una vez al
nodo final (\in[eq:mono_llegada]). Además, no vuelve al depósito
(\in[eq:mono_no_vuelve]) ni sale del nodo final (\in[eq:mono_no_sale]).
\startplaceformula
\startformula
	\startalign[n=3, align={right,left,left}]
		\NC \NC \sum_{(i,j) \in \delta^{+}(v_{\text{start}})} z_{ijb} = 1 \NC \qquad \forall b \in \mathcal{B} \NR[eq:mono_salida]
		\NC \NC \sum_{(i,j) \in \delta^{-}(v_{\text{end}})} z_{ijb} = 1 \NC \qquad \forall b \in \mathcal{B} \NR[eq:mono_llegada]
		\NC \NC \sum_{(i,j) \in \delta^{-}(v_{\text{start}})} z_{ijb} = 0 \NC \qquad \forall b \in \mathcal{B} \NR[eq:mono_no_vuelve]
		\NC \NC \sum_{(i,j) \in \delta^{+}(v_{\text{end}})} z_{ijb} = 0 \NC \qquad \forall b \in \mathcal{B} \NR[eq:mono_no_sale]
	\stopalign
\stopformula
\stopplaceformula

{\bf Visitas y continuidad.}
A cada ubicación requerida se entra (\in[eq:mono_entra]) y se sale
(\in[eq:mono_sale]). En las ubicaciones intermedias, lo que entra sale
(\in[eq:mono_flujo]). A cada ubicación se entra (\in[eq:mono_entra_max]) y
se sale (\in[eq:mono_sale_max]) como máximo una vez, de modo que ninguna se
repite.
% TODO: Acá estaría bueno explicar cada una de estas restricciones, si no es redundante
\startplaceformula
\startformula
	\startalign[n=3, align={right,left,left}]
		\NC \NC \sum_{(i,j) \in \delta^{-}(j)} z_{ijb} \geq a_{bj} \NC \qquad \forall b \in \mathcal{B},\; j \in \mathcal{V}^{\text{int}} \NR[eq:mono_entra]
		\NC \NC \sum_{(j,i) \in \delta^{+}(j)} z_{jib} \geq a_{bj} \NC \qquad \forall b \in \mathcal{B},\; j \in \mathcal{V}^{\text{int}} \NR[eq:mono_sale]
		\NC \NC \sum_{(g,i) \in \delta^{-}(i)} z_{gib} = \sum_{(i,j) \in \delta^{+}(i)} z_{ijb} \NC \qquad \forall b \in \mathcal{B},\; i \in \mathcal{V}^{\text{int}} \NR[eq:mono_flujo]
		\NC \NC \sum_{(j,i) \in \delta^{-}(i)} z_{jib} \leq 1 \NC \qquad \forall b \in \mathcal{B},\; i \in \mathcal{V} \setminus \{v_{\text{start}}\} \NR[eq:mono_entra_max]
		\NC \NC \sum_{(i,j) \in \delta^{+}(i)} z_{ijb} \leq 1 \NC \qquad \forall b \in \mathcal{B},\; i \in \mathcal{V} \setminus \{v_{\text{end}}\} \NR[eq:mono_sale_max]
	\stopalign
\stopformula
\stopplaceformula

{\bf Eliminación de subtours.}
Las restricciones (\in[eq:mono_salida])--(\in[eq:mono_sale_max]) todavía
admiten soluciones formadas por un camino del depósito al nodo final más un
ciclo disjunto que visita algunas ubicaciones requeridas sin conectarse con el
depósito. Para eliminarlos se utiliza una variante de la formulación de flujo
de Gavish y Graves (Ecuación~\in[eq:gg]). Un flujo $q$ sale del depósito con
$L_b$ unidades, una por cada ubicación requerida, y cada ubicación requerida
consume una. Como el flujo solo circula por arcos usados
(\in[eq:mono_gg_cota]) y se origina únicamente en el depósito
(\in[eq:mono_gg_origen]), una ubicación requerida solo puede recibir su
unidad (\in[eq:mono_gg_balance]) si está conectada al depósito por la ruta.
\startplaceformula
\startformula
	\startalign[n=3, align={right,left,left}]
		\NC \NC q_{ijb} \leq L_b \, z_{ijb} \NC \qquad \forall b \in \mathcal{B},\; (i,j) \in \mathcal{E} \NR[eq:mono_gg_cota]
		\NC \NC \sum_{(i,j) \in \delta^{+}(v_{\text{start}})} q_{ijb} = L_b \NC \qquad \forall b \in \mathcal{B} \NR[eq:mono_gg_origen]
		\NC \NC \sum_{(g,i) \in \delta^{-}(i)} q_{gib} - \sum_{(i,j) \in \delta^{+}(i)} q_{ijb} = a_{bi} \NC \qquad \forall b \in \mathcal{B},\; i \in \mathcal{V}^{\text{int}} \NR[eq:mono_gg_balance]
	\stopalign
\stopformula
\stopplaceformula

El flujo, sin embargo, solo obliga a conectar las ubicaciones requeridas, de
modo que un ciclo disjunto que pase únicamente por ubicaciones no requeridas
sigue siendo factible. Como ese ciclo solo agrega tiempo de ruta, nunca mejora el objetivo.
En el modelo monolítico podría aparecer únicamente en una solución óptima
alternativa, cuando el tiempo de ruta no influye en el objetivo, y la Etapa~1
del modelo Split, que minimiza el tiempo de ruta, no lo produce.

{\bf Tiempo de ruta.}
El tiempo de ruta es la distancia recorrida multiplicada por
$T^{\text{travel}}$, más el tiempo de servicio de cada ubicación intermedia
por la que pasa la ruta (una ubicación se cuenta cuando la ruta sale de ella):
\startplaceformula[reference=eq:mono_ruta]
\startformula
	r_b = T^{\text{travel}} \sum_{(i,j) \in \mathcal{E}} D_{ij} \, z_{ijb}
	\;+\; \sum_{i \in \mathcal{V}^{\text{int}}} S_i \sum_{(i,j) \in \delta^{+}(i)} z_{ijb}
	\qquad \forall b \in \mathcal{B}
\stopformula
\stopplaceformula
El tiempo de servicio se cobra en toda ubicación intermedia que la ruta
atraviesa, se recoja algo en ella o no.\footnote{En la implementación, la
suma del traslado excluye los arcos que entran al depósito y los que salen del
nodo final. Por (\in[eq:mono_no_vuelve]) y (\in[eq:mono_no_sale]) esos arcos
no se usan, así que las dos expresiones coinciden.}

\stopsubsection

\startsubsection[title={Asignación, secuencia y fechas de entrega}, reference=sec:mono_asignacion]

% TODO: Acá se ve claro lo de M, quizás no hace falta tanta explicación previa
{\bf Asignación y tiempo de ejecución.} Cada lote se asigna a exactamente un
{\em picker} (\in[eq:mono_asignacion]). Si el {\em picker} $p$ ejecuta el
lote $b$, las restricciones (\in[eq:mono_u_inf]) y (\in[eq:mono_u_sup]), de
tipo M grande (Ecuación~\in[eq:bigm]), fijan $u_{pb}$ exactamente en su
tiempo de ejecución $F_b + r_b / V_p$ (Ecuación~\in[eq:texec]); si no lo
ejecuta, quedan desactivadas, para lo cual $M$ debe ser al menos el mayor
tiempo de ejecución posible de un lote. El lote
no puede terminar antes de lo que tarda en ejecutarse (\in[eq:mono_piso]), y
su instante de finalización es el de la agenda del {\em picker} que lo
ejecuta (\in[eq:mono_cb]).
\startplaceformula
\startformula
	\startalign[n=3, align={right,left,left}]
		\NC \NC \sum_{p \in \mathcal{P}} x_{pb} = 1 \NC \qquad \forall b \in \mathcal{B} \NR[eq:mono_asignacion]
		\NC \NC u_{pb} \geq F_b + \frac{r_b}{V_p} - M (1 - x_{pb}) \NC \qquad \forall p \in \mathcal{P},\; b \in \mathcal{B} \NR[eq:mono_u_inf]
		\NC \NC u_{pb} \leq F_b + \frac{r_b}{V_p} + M (1 - x_{pb}) \NC \qquad \forall p \in \mathcal{P},\; b \in \mathcal{B} \NR[eq:mono_u_sup]
		\NC \NC c_{pb} \geq u_{pb} \NC \qquad \forall p \in \mathcal{P},\; b \in \mathcal{B} \NR[eq:mono_piso]
		\NC \NC c_b \geq c_{pb} - M (1 - x_{pb}) \NC \qquad \forall p \in \mathcal{P},\; b \in \mathcal{B} \NR[eq:mono_cb]
	\stopalign
\stopformula
\stopplaceformula

{\bf Cadena de frío.} No necesita una restricción aparte, ya que está contenida
en la cota superior de la variable, $u_{pb} \leq T^{\text{cold}} + \varepsilon$.
Si $x_{pb} = 1$, la restricción (\in[eq:mono_u_inf]) obliga a que
$u_{pb} \geq F_b + r_b / V_p$, por lo que el {\em picker} $p$ solo puede
ejecutar el lote $b$ si $F_b + r_b / V_p \leq T^{\text{cold}} + \varepsilon$.
Si $x_{pb} = 0$, $u_{pb}$ puede valer 0 y la cota no restringe nada. Como el
límite se aplica a cada lote y no a la jornada del {\em picker}, un mismo
{\em picker} puede ejecutar varios lotes seguidos aunque el último termine
después de $T^{\text{cold}}$.

{\bf Secuencia sin solapamiento.} Un {\em picker} ejecuta un lote por vez.
Para cada {\em picker} $p$ y cada par de lotes $b < k$:
\startplaceformula
\startformula
	\startalign[n=3, align={right,left,left}]
		\NC \NC c_{pk} \geq c_{pb} + F_k + \frac{r_k}{V_p} - M^{\text{seq}} (3 - x_{pb} - x_{pk} - y_{bk}) \NC \quad \forall p \in \mathcal{P},\; b, k \in \mathcal{B},\; b < k \NR[eq:mono_seq1]
		\NC \NC c_{pb} \geq c_{pk} + F_b + \frac{r_b}{V_p} - M^{\text{seq}} (2 - x_{pb} - x_{pk} + y_{bk}) \NC \quad \forall p \in \mathcal{P},\; b, k \in \mathcal{B},\; b < k \NR[eq:mono_seq2]
	\stopalign
\stopformula
\stopplaceformula
Cada una de estas restricciones se activa solo cuando tres variables
binarias toman un valor dado, con la técnica de la
Sección~\in[sec:marco_mip]. La primera se activa cuando
$x_{pb} = x_{pk} = y_{bk} = 1$, ya que solo entonces el factor
$3 - x_{pb} - x_{pk} - y_{bk}$ vale 0. La segunda usa $1 - y_{bk}$ en lugar de $y_{bk}$, porque su
condición es $y_{bk} = 0$. Así, si $p$ ejecuta ambos lotes, solo está activa
la restricción que corresponde al orden elegido, y el segundo lote termina al
menos su tiempo de ejecución después del primero. Si $p$ no ejecuta alguno de
los dos, ambas quedan desactivadas.
% TODO: Ver si esto va acá o va en el marco teórico
Para que la restricción inactiva no limite la solución, $M^{\text{seq}}$ debe
superar la mayor separación posible entre las finalizaciones de dos lotes de
un mismo {\em picker}, sumada a la duración de un lote. La cota de la
Sección~\in[sec:notacion] cumple esta condición cuando cada lote termina en
cuanto completa su ejecución; la Sección~\in[sec:caracteristicas_formulacion]
analiza el caso en que esto no ocurre.

{\bf Adelanto y demora.} Cada orden termina cuando termina su lote. Su demora es el tiempo en que se
excede de la fecha de entrega, y su adelanto, el tiempo que le falta para
alcanzarla:
\startplaceformula
\startformula
	\startalign[n=3, align={right,left,left}]
		\NC \NC ta_n \geq c_{pb} - T^{\text{due}}_n - M (1 - x_{pb}) \NC \qquad \forall p \in \mathcal{P},\; b \in \mathcal{B},\; n \in \mathcal{N}_b \NR[eq:mono_demora]
		\NC \NC ea_n \geq T^{\text{due}}_n - c_{pb} - M (1 - x_{pb}) \NC \qquad \forall p \in \mathcal{P},\; b \in \mathcal{B},\; n \in \mathcal{N}_b \NR[eq:mono_adelanto]
	\stopalign
\stopformula
\stopplaceformula
Si $x_{pb} = 0$, las dos restricciones quedan desactivadas, siempre que $M$
supere la fecha de entrega más lejana y el horizonte de
planificación.\footnote{En la implementación, estas restricciones se escriben
con el término $M (2 - x_{pb} - B_{nb})$ y se generan solo para los pares con
$B_{nb} = 1$, lo que es equivalente.} Como $ta_n, ea_n \geq 0$ y el objetivo
los minimiza, en el óptimo se cumple que
$ta_n = \max(0,\, c_b - T^{\text{due}}_n)$ y
$ea_n = \max(0,\, T^{\text{due}}_n - c_b)$.

\stopsubsection
\stopsection

%----------------------------------------------------------------------------------------
\startsection[title={Modelo en Dos Etapas (Split)}, reference=sec:modelo_split]

El modelo Split resuelve el mismo problema mediante dos \ac{mip} que se
ejecutan uno después del otro:
\startitemize[n]
	\item {\bf Etapa 1.} Calcula la ruta más corta de cada lote.
	\item {\bf Etapa 2.} Con esos tiempos de ruta como datos, decide la
	asignación, la secuencia y el cumplimiento de la cadena de frío.
\stopitemize

La idea de separar el ruteo del resto de las decisiones se parece a la del
modelo simplificado de Coruzzolo et al. (Sección~\in[sec:trabajos_relacionados]),
con dos diferencias. Las rutas se calculan de forma exacta, y no con la
heurística del vecino más cercano, y la descomposición alcanza el mismo valor
óptimo que el modelo monolítico, como se muestra a continuación.

La descomposición se apoya en dos propiedades. La primera es que la ruta de
un lote no depende de qué {\em picker} lo ejecuta ni del orden entre lotes,
por los supuestos de velocidad y de ausencia de congestión. La segunda es que
reemplazar la ruta de un lote por una más corta nunca empeora el objetivo.
Para verlo, considérese una solución óptima del modelo monolítico y
reemplácese la ruta de cada lote por su ruta mínima. El tiempo de ejecución
$F_b + r_b / V_p$ de cada lote no aumenta, de modo que la cadena de frío se
sigue cumpliendo. Además, como las restricciones que vinculan la ejecución con
la finalización, (\in[eq:mono_piso]) y
(\in[eq:mono_seq1])--(\in[eq:mono_seq2]), son desigualdades, los mismos
instantes de finalización $c_{pb}$, la misma asignación $x_{pb}$ y la misma
secuencia $y_{bk}$ siguen siendo factibles. El objetivo depende solo de las
finalizaciones y de los adelantos y demoras, así que no cambia. Existe, por lo
tanto, una solución óptima con rutas mínimas, que es la que busca el modelo
Split, y ambos modelos alcanzan el mismo valor óptimo.

El argumento usa que un lote puede figurar como terminado más tarde de lo que
termina realmente, propiedad que se conoce como espera voluntaria
(Sección~\in[sec:espera_voluntaria]). Si las finalizaciones quedaran
determinadas por igualdades, acortar una ruta adelantaría la finalización del
lote y podría aumentar su adelanto, y la descomposición dejaría de ser exacta
en general.

% TODO: Esto necesita más explicación - tenemos que dejar claro por qué esto
% es así. Notar que pasa porque el modelo monolítico permite rutas no mínimas,
% pero el modelo Split siempre busca la ruta mínima en la Etapa 1. La pregunta es,
% hay una ruta no mínima que genere un mejor valor objetivo que la ruta mínima? 
% La respuesta es no, porque si reemplazamos la ruta no mínima por la mínima, 
% el tiempo de ejecución del lote no aumenta, y por lo tanto el valor objetivo 
% no empeora. Por lo tanto, siempre podemos encontrar una solución óptima con rutas mínimas.
% Otro tema es: por qué el modelo monolítico puede devolver rutas no mínimas? 
% Porque no tiene un objetivo que minimice la ruta, solo minimiza el tiempo de 
% finalización y las penalizaciones. Por lo tanto, puede haber múltiples rutas 
% que den el mismo tiempo de finalización y penalizaciones, y algunas de esas 
% rutas pueden no ser mínimas. En ese caso, el valor objetivo sería el mismo, 
% pero la ruta no sería mínima. Esto no es un problema del modelo en sí, es una 
% diferencia de enfoque de los modelos. Esto se puede charlar acá o ampliar en otra
% sección, porque es una consecuencia interesante. También se pueden tomar 
% resultados de los experimentos para respaldar estas afirmaciones (ambas, la que
% dice que el modelo monolítico puede devolver rutas no mínimas y la que dice que
% por más que la ruta no sea mínima, el valor objetivo no empeora). 
% Esto se puede poner en la sección de experimentos, donde se comparan 
% los resultados de ambos modelos.
El conjunto de soluciones óptimas, en cambio, no es el mismo. El modelo
monolítico puede devolver rutas no mínimas que alcanzan el mismo valor
objetivo, mientras que el modelo Split siempre devuelve rutas mínimas.

\startsubsection[title={Etapa 1: ruteo de cada lote}, reference=sec:split_etapa1]

{\bf Variables.} $z_{ijb}$, $q_{ijb}$ y $r_b$, con el mismo
significado que en el modelo monolítico.

{\bf Función objetivo.} Minimizar el tiempo total de ruta:
\startplaceformula[reference=eq:split_objetivo1]
\startformula
	\min \quad \sum_{b \in \mathcal{B}} r_b
\stopformula
\stopplaceformula
% TODO: acá hablamos de "resolver en paralelo". Esto no está mal, porque es uno
% de los motivos justamente por los que dividimos el modelo en dos etapas. No es
% la única optimización que se puede hacer (y se probaron más). El tema es que no
% sé si este capítulo sea el lugar para hablarlo.
Como los lotes no comparten ninguna restricción, minimizar la suma equivale a
minimizar cada $r_b$ por separado, por lo que en la práctica cada lote se
resuelve como un \ac{mip} independiente. Al no compartir variables ni
restricciones, esos modelos también pueden resolverse en paralelo.

{\bf Restricciones.} Las mismas restricciones de ruteo del modelo
monolítico, (\in[eq:mono_salida])--(\in[eq:mono_ruta]). Esta etapa no
utiliza {\em pickers}, pesos del objetivo ni tiempos de preparación o de
recolección.

{\bf Resultado.} El tiempo de ruta mínimo de cada lote,
$R_b = r_b^{*}$, que pasa a la Etapa~2 como dato.

\stopsubsection

\startsubsection[title={Etapa 2: asignación, secuencia y cadena de frío}, reference=sec:split_etapa2]

{\bf Datos provenientes de la Etapa 1.}
Con $R_b$ fijo, el tiempo de ejecución de cada par ({\em picker}, lote) es
un valor conocido antes de resolver, y la cadena de frío puede verificarse de
antemano:
\startformula
	\startalign[n=2, align={right,left}]
		\NC T^{\text{exec}}_{pb} \EQ F_b + \frac{R_b}{V_p}, \NR
		\NC \mathcal{PB}^{\text{elig}} \EQ \left\{ (p,b) \in \mathcal{P} \times \mathcal{B} : T^{\text{exec}}_{pb} \leq T^{\text{cold}} + \varepsilon \right\}. \NR
	\stopalign
\stopformula
$\mathcal{PB}^{\text{elig}}$ es el conjunto de pares elegibles, es decir,
aquellos en que el {\em picker} puede ejecutar el lote sin violar la cadena
de frío. Las variables $x_{pb}$ y $c_{pb}$ se definen solo para esos pares,
de modo que la cadena de frío se cumple por construcción y no necesita
ninguna restricción. Si algún lote no tiene ningún {\em picker} elegible, la
restricción de asignación (\in[eq:split_asignacion]) no puede cumplirse y la
instancia es infactible.

{\bf Variables.}
\startformula
	\startalign[n=2, align={right,left}]
		\NC x_{pb} \in \{0,1\} \NC \quad \text{1 si el lote $b$ se asigna al {\em picker} $p$, para $(p,b) \in \mathcal{PB}^{\text{elig}}$} \NR
		\NC c_{pb} \geq 0 \NC \quad \text{instante en que termina el lote $b$ en la agenda de $p$, para $(p,b) \in \mathcal{PB}^{\text{elig}}$} \NR
		\NC c_b \geq 0 \NC \quad \text{instante en que termina el lote $b$} \NR
		\NC y_{bk} \in \{0,1\} \NC \quad \text{para $b < k$: 1 si $b$ se ejecuta antes que $k$} \NR
		\NC ta_n \geq 0,\; ea_n \geq 0 \NC \quad \text{demora y adelanto de la orden $n$} \NR
	\stopalign
\stopformula
La variable $u_{pb}$ del modelo monolítico ya no es necesaria, porque el
tiempo de ejecución es un dato.

{\bf Función objetivo.} La misma que la del modelo monolítico:
\startplaceformula[reference=eq:split_objetivo2]
\startformula
	\min \quad W^{c} \sum_{b \in \mathcal{B}} c_b \;+\; W^{\text{pen}} \left( W^{e} \sum_{n \in \mathcal{N}} ea_n + W^{t} \sum_{n \in \mathcal{N}} ta_n \right)
\stopformula
\stopplaceformula

{\bf Asignación y finalización.}
Cada lote se asigna a exactamente un {\em picker} elegible
(\in[eq:split_asignacion]). Si el {\em picker} $p$ ejecuta el lote $b$, el
lote no termina antes de su tiempo de ejecución (\in[eq:split_piso]); si no
lo ejecuta, el lado derecho es 0 y la restricción no restringe. El instante
de finalización del lote es el de la agenda del {\em picker} que lo ejecuta
(\in[eq:split_cb]).
\startplaceformula
\startformula
	\startalign[n=3, align={right,left,left}]
		\NC \NC \sum_{p \,:\, (p,b) \in \mathcal{PB}^{\text{elig}}} x_{pb} = 1 \NC \qquad \forall b \in \mathcal{B} \NR[eq:split_asignacion]
		\NC \NC c_{pb} \geq T^{\text{exec}}_{pb} \, x_{pb} \NC \qquad \forall (p,b) \in \mathcal{PB}^{\text{elig}} \NR[eq:split_piso]
		\NC \NC c_b \geq c_{pb} - M (1 - x_{pb}) \NC \qquad \forall (p,b) \in \mathcal{PB}^{\text{elig}} \NR[eq:split_cb]
	\stopalign
\stopformula
\stopplaceformula
La restricción (\in[eq:split_piso]) reemplaza a
(\in[eq:mono_u_inf])--(\in[eq:mono_piso]) del modelo monolítico. Como
$T^{\text{exec}}_{pb}$ es un número, el producto
$T^{\text{exec}}_{pb} \, x_{pb}$ es lineal y no hace falta una restricción
de tipo M grande.

{\bf Secuencia sin solapamiento.}
Son las mismas restricciones (\in[eq:mono_seq1])--(\in[eq:mono_seq2]) del
modelo monolítico, con $T^{\text{exec}}_{pb}$ como dato. Se generan para
cada {\em picker} $p$ y cada par de lotes $b < k$ elegibles para $p$:
\startplaceformula
\startformula
	\startalign[n=3, align={right,left,left}]
		\NC \NC c_{pk} \geq c_{pb} + T^{\text{exec}}_{pk} - M^{\text{seq}} (3 - x_{pb} - x_{pk} - y_{bk}) \NC \quad \forall (p,b), (p,k) \in \mathcal{PB}^{\text{elig}},\; b < k \NR[eq:split_seq1]
		\NC \NC c_{pb} \geq c_{pk} + T^{\text{exec}}_{pb} - M^{\text{seq}} (2 - x_{pb} - x_{pk} + y_{bk}) \NC \quad \forall (p,b), (p,k) \in \mathcal{PB}^{\text{elig}},\; b < k \NR[eq:split_seq2]
	\stopalign
\stopformula
\stopplaceformula
Como en el modelo monolítico, solo una de las dos está activa, según qué
lote va primero, y ambas se desactivan si $p$ no ejecuta los dos lotes.

{\bf Adelanto y demora.}
Son las mismas restricciones (\in[eq:mono_demora])--(\in[eq:mono_adelanto])
del modelo monolítico, restringidas a los pares elegibles:
\startplaceformula
\startformula
	\startalign[n=3, align={right,left,left}]
		\NC \NC ta_n \geq c_{pb} - T^{\text{due}}_n - M (1 - x_{pb}) \NC \qquad \forall (p,b) \in \mathcal{PB}^{\text{elig}},\; n \in \mathcal{N}_b \NR[eq:split_demora]
		\NC \NC ea_n \geq T^{\text{due}}_n - c_{pb} - M (1 - x_{pb}) \NC \qquad \forall (p,b) \in \mathcal{PB}^{\text{elig}},\; n \in \mathcal{N}_b \NR[eq:split_adelanto]
	\stopalign
\stopformula
\stopplaceformula

\stopsubsection

\startsubsection[title={Resolución}, reference=sec:split_resolucion]

La Figura~\in[fig:split_flujo] muestra cómo se resuelve una instancia con el
modelo Split.
% TODO: revisar la validación de los datos (antes, el ítem 1 de una lista).
Primero se validan los datos. Se verifica que cada orden pertenezca a un
único lote, que ningún lote supere la capacidad del carro y que todas las
velocidades sean positivas, ya que una velocidad nula dejaría indefinido el
tiempo de ejecución y una negativa lo haría negativo. Si alguna de estas
condiciones no se cumple, la instancia se informa como infactible sin
construir ningún modelo.

Después se resuelve la Etapa~1, que consiste en un \ac{mip} de ruteo por
lote, independientes entre sí (Sección~\in[sec:split_etapa1]). Cada uno
devuelve el tiempo de ruta mínimo $R_b$ de su lote. Con esos valores se
calculan los tiempos de ejecución $T^{\text{exec}}_{pb}$ y el conjunto de
pares elegibles $\mathcal{PB}^{\text{elig}}$. Si algún lote no tiene ningún
{\em picker} elegible, la instancia se informa como infactible sin resolver
la Etapa~2. En caso contrario, se resuelve la Etapa~2, que decide la
asignación y la secuencia con los tiempos de ejecución como datos y completa
la solución.

El modelo monolítico comparte la validación de los datos, pero después
resuelve un único \ac{mip} que decide a la vez las rutas, la asignación y la
secuencia. En ese caso, si la cadena de frío vuelve infactible la instancia,
lo detecta el resolvedor.

\startplacefigure[title={Flujo de resolución del modelo Split. Las flechas rojas marcan los casos en que la instancia se informa como infactible antes de resolver la Etapa~2.}, reference=fig:split_flujo]
\startMPcode
numeric w, h, dy, hs; w := 7.6cm; h := 0.9cm; dy := 1.5cm; hs := 1.8cm;
color borde, relleno, rojo;
borde := \MPcolor{darkblue}; relleno := \MPcolor{darkgray}; rojo := \MPcolor{redline};
def caja(expr cy, hh, txt) =
  fill unitsquare xyscaled (w, hh) shifted (-w/2, cy - hh/2) withcolor relleno;
  draw unitsquare xyscaled (w, hh) shifted (-w/2, cy - hh/2) withpen pencircle scaled .8pt withcolor borde;
  draw textext(txt) shifted (0, cy);
enddef;
def flecha(expr ya, yb) =
  drawarrow (0, ya) -- (0, yb) withpen pencircle scaled .8pt;
enddef;
def salida(expr cy, txt) =
  drawarrow (w/2, cy) -- (w/2 + 0.9cm, cy) withpen pencircle scaled .8pt withcolor rojo;
  label.rt(textext(txt), (w/2 + 0.95cm, cy)) withcolor rojo;
enddef;
numeric y[];
y0 := 0; y1 := y0 - dy; y2 := y1 - h/2 - 0.6cm - hs/2; y3 := y2 - hs/2 - 0.9cm - h/2;
y4 := y3 - dy; y5 := y4 - dy;
caja(y0, h, "Datos de la instancia");
flecha(y0 - h/2, y1 + h/2);
caja(y1, h, "Validación de los datos");
salida(y1, "\framed[frame=off, align=flushleft, offset=0pt]{\tfx infactible:\crlf datos inválidos}");
flecha(y1 - h/2, y2 + hs/2);
caja(y2, hs, "");
draw textext("Etapa 1: un MIP de ruteo por cada lote") shifted (0, y2 + 0.45cm);
numeric xs[]; xs1 := -2.55cm; xs2 := -0.85cm; xs3 := 0.85cm; xs4 := 2.55cm;
string ls[]; ls1 := "\tfx lote 1"; ls2 := "\tfx lote 2"; ls3 := "$\cdots$"; ls4 := "\tfx lote $|\mathcal{B}|$";
for i = 1 upto 4:
  fill unitsquare xyscaled (1.4cm, 0.6cm) shifted (xs[i] - 0.7cm, y2 - 0.6cm) withcolor white;
  draw unitsquare xyscaled (1.4cm, 0.6cm) shifted (xs[i] - 0.7cm, y2 - 0.6cm) withpen pencircle scaled .5pt withcolor borde;
  draw textext(ls[i]) shifted (xs[i], y2 - 0.3cm);
endfor;
flecha(y2 - hs/2, y3 + h/2);
label.rt(textext("$R_b$"), (0.1cm, (y2 - hs/2 + y3 + h/2)/2));
caja(y3, h, "Elegibilidad: $T^{\text{exec}}_{pb}$ y $\mathcal{PB}^{\text{elig}}$");
salida(y3, "\framed[frame=off, align=flushleft, offset=0pt]{\tfx infactible: algún lote\crlf sin {\em picker} elegible}");
flecha(y3 - h/2, y4 + h/2);
caja(y4, h, "Etapa 2: un MIP de asignación y secuencia");
flecha(y4 - h/2, y5 + h/2);
caja(y5, h, "Solución: rutas, asignación y secuencia");
\stopMPcode
\stopplacefigure

\stopsubsection
\stopsection

%----------------------------------------------------------------------------------------
\startsection[title={Espera voluntaria}, reference=sec:espera_voluntaria]

En ambos modelos, un lote puede terminar más tarde de lo que permite su
ejecución, porque el {\em picker} puede demorar su inicio. Esto le permite
acercar su instante de finalización a la fecha de entrega de sus órdenes y
pagar menos penalización por adelanto. A este comportamiento se lo denomina
{\em espera voluntaria}, y es la forma que toma en estos modelos el tiempo
ocioso insertado de los problemas con adelanto y demora
(Sección~\in[sec:marco_adelanto_demora]). Como las metaheurísticas del
Capítulo~\in[Capitulo_4] no generan esta espera, hay que tenerla en cuenta al
comparar sus resultados con los de los modelos lineales.

\startsubsection[title={Por qué ocurre}]

% TODO: acá es importante aclarar por qué usamos desigualdades. Puede verse como
% que las desigualdades son incorrectas, pero son necesarias para modelar lo que
% queremos modelar.
Las restricciones que vinculan la ejecución de un lote con su instante de
finalización son desigualdades, no igualdades. En el modelo monolítico,
(\in[eq:mono_u_inf])--(\in[eq:mono_u_sup]) fijan $u_{pb}$ en la duración
real del lote, pero (\in[eq:mono_piso]) solo exige que
$c_{pb} \geq u_{pb}$. Lo mismo ocurre al encadenar lotes del mismo
{\em picker} en (\in[eq:mono_seq1])--(\in[eq:mono_seq2]), que solo exigen
que $c_{pk}$ sea al menos $c_{pb}$ más la duración de $k$. En el modelo
Split sucede lo mismo con (\in[eq:split_piso]) y
(\in[eq:split_seq1])--(\in[eq:split_seq2]). Nada impide, entonces, que
$c_{pb}$ quede por encima de su cota inferior, ya sea al final de la
secuencia de un {\em picker} o entre dos lotes consecutivos. El modelo de
referencia tiene la misma propiedad, porque también encadena las
finalizaciones de los lotes de un operario con desigualdades
\cite[coruzzolo2023orderpicking].

Por otro lado, la función objetivo (\in[eq:mono_objetivo]) penaliza también
el adelanto, de modo que, si un lote termina antes de la fecha de entrega de
sus órdenes, al resolvedor puede convenirle declarar un $c_b$ mayor. Esta misma
holgura es la que garantiza que el modelo Split alcance el óptimo del
monolítico (Sección~\in[sec:modelo_split]).

% TODO: creo que nos hace falta un ejemplo de literatura. No es un concepto que
% nos inventamos. También estaría bueno mencionar que el modelo original de 
% \cite{baker1981sequencing} tiene espera voluntaria, y justificar por qué.

\stopsubsection

\startsubsection[title={Cuándo esperar es conveniente}]

Si se demora en $\delta > 0$ unidades el instante de finalización de un
lote, cada término del objetivo cambia como se resume en el
Cuadro~\in[tab:espera_balance].

\startplacetable[reference=tab:espera_balance, location=here, title={Efecto en la función objetivo de demorar en $\delta$ el instante de finalización de un lote.}]
	\bTABLE[split=yes, align=middle, frame=on, width=fit]
	\setupTABLE[row][1][background=color, backgroundcolor=darkgray, style=bold]
	\setupTABLE[column][1][width=7cm, align=flushleft]
	\setupTABLE[column][2][width=6cm, align=flushleft]
	\bTR \bTH Término \eTH \bTH Efecto \eTH \eTR
	\bTR \bTD Finalización del lote \eTD \bTD $+\, W^{c} \, \delta$ \eTD \eTR
	\bTR \bTD Adelanto de cada orden del lote que todavía está antes de su fecha de entrega \eTD \bTD $-\, W^{\text{pen}} W^{e} \, \delta$ por orden \eTD \eTR
	\bTR \bTD Demora de cada orden del lote que ya está atrasada \eTD \bTD $+\, W^{\text{pen}} W^{t} \, \delta$ por orden \eTD \eTR
	\bTR \bTD Lotes posteriores del mismo {\em picker}, si no tienen holgura \eTD \bTD se desplazan también y pagan su propio costo \eTD \eTR
	\eTABLE
\stopplacetable

Sean $n^{e}_b$ y $n^{t}_b$ la cantidad de órdenes del lote $b$ que están
adelantadas y atrasadas, respectivamente, y $\Delta^{\text{suc}}_b$ el costo
marginal de desplazar a los lotes posteriores del mismo {\em picker}. La
espera reduce el objetivo cuando
\startplaceformula[reference=eq:espera_condicion]
\startformula
	W^{\text{pen}} W^{e} \, n^{e}_b \;>\; W^{c} + W^{\text{pen}} W^{t} \, n^{t}_b + \Delta^{\text{suc}}_b.
\stopformula
\stopplaceformula
Con $W^{\text{pen}} = W^{e} = 1$, y para un lote sin órdenes atrasadas ni
sucesores que desplazar, la condición se reduce a $n^{e}_b > W^{c}$. La
espera nunca conviene si el adelanto no tiene peso en el objetivo
($W^{\text{pen}} W^{e} = 0$), y tampoco en ningún lote si
$W^{c} \geq W^{\text{pen}} W^{e} \max_{b \in \mathcal{B}} |\mathcal{N}_b|$.

Una consecuencia de esta condición es que, cuando la espera conviene, la
proporción de órdenes que llegan a tiempo queda determinada por los pesos del
objetivo y por el tamaño del lote (Anexo~\in[Anexo_Modelos]).

El costo en función de $c_b$ es lineal por tramos, y cada fecha de entrega
es un punto de quiebre ya que, al superarla, una orden deja de aportar adelanto y
empieza a aportar demora. Por eso, cuando conviene esperar, el instante de
finalización óptimo coincide con la fecha de entrega de alguna orden, de
forma análoga a una mediana ponderada.

% TODO: gráfico, o algo que ayude a visualizar el ejemplo
{\bf Ejemplo.} Sea un lote con tres órdenes cuya fecha de entrega es
20, que termina realmente en $t = 10$, con todos los pesos iguales a 1. Si se
reporta su instante real de finalización, su aporte al objetivo es
$1 \cdot 10 + 3 \cdot 10 = 40$ (finalización más adelanto de las tres
órdenes). Si se lo declara terminado en $t = 20$, el aporte es
$1 \cdot 20 + 0 = 20$. Más allá de $t = 20$, cada unidad adicional cuesta
$1 + 3 = 4$ (finalización más demora de las tres órdenes), por lo que el
resolvedor se detiene exactamente en la fecha de entrega.

% TODO: Podemos mencionar que lo vamos a ver en los experimentos... y poner el caso
% que tenemos en los experimentos. Quizás podemos armar uno que otro test case
% funcional y ejecutarlo en ambos modelos, así vemos claramente esto.

\stopsubsection
\stopsection

%----------------------------------------------------------------------------------------
\startsection[title={Diferencias entre las formulaciones}, reference=sec:diferencias_modelos]

El Cuadro~\in[tab:diferencias_modelos] resume las diferencias entre el
modelo monolítico y el modelo Split.

% TODO: No se si me convence esto. Quizás en formato lista es mejor.
\startplacetable[reference=tab:diferencias_modelos, location=here, title={Diferencias entre el modelo monolítico y el modelo Split.}]
	\bTABLE[split=yes, align=middle, frame=on, width=fit]
	\setupTABLE[row][1][background=color, backgroundcolor=darkgray, style=bold]
	\setupTABLE[column][1][width=3.6cm, align=flushleft]
	\setupTABLE[column][2][width=4.8cm, align=flushleft]
	\setupTABLE[column][3][width=4.8cm, align=flushleft]
	\bTR \bTH Aspecto \eTH \bTH Monolítico \eTH \bTH Split \eTH \eTR
	\bTR \bTD Cantidad de \ac{mip} \eTD \bTD 1 \eTD \bTD $1 + |\mathcal{B}|$: uno de ruteo por lote (Etapa~1), independientes entre sí, y uno de asignación y secuencia (Etapa~2) \eTD \eTR
	\bTR \bTD Tiempos de ruta y de ejecución \eTD \bTD Variables de decisión ($r_b$ y $u_{pb}$, con M grande) \eTD \bTD Datos ($R_b$, de la Etapa~1, y $T^{\text{exec}}_{pb}$) \eTD \eTR
	\bTR \bTD Cadena de frío \eTD \bTD Cota superior de $u_{pb}$; si la vuelve infactible, lo detecta el resolvedor \eTD \bTD Pares elegibles $\mathcal{PB}^{\text{elig}}$; la infactibilidad se detecta antes de la Etapa~2 \eTD \eTR
	\bTR \bTD Ruta mínima garantizada \eTD \bTD No (puede elegir una ruta no mínima) \eTD \bTD Sí (la Etapa 1 minimiza $r_b$) \eTD \eTR
	\bTR \bTD Tamaño \eTD \bTD $2 |\mathcal{E}| |\mathcal{B}|$ variables de ruteo en el mismo \ac{mip} que la asignación y la secuencia \eTD \bTD $2 |\mathcal{E}| + 1$ variables por cada \ac{mip} de la Etapa~1; la Etapa~2 no depende del grafo \eTD \eTR
	
	\eTABLE
\stopplacetable

\stopsection

%----------------------------------------------------------------------------------------
\startsection[title={Características de la formulación}, reference=sec:caracteristicas_formulacion]

Las dos formulaciones comparten algunas características que conviene tener
presentes al interpretar sus soluciones.

{\bf Camino simple.} Las restricciones (\in[eq:mono_entra_max]) y
(\in[eq:mono_sale_max]) impiden que una ruta pase dos veces por la misma
ubicación. Es una característica que se hereda del modelo de referencia y que resulta
adecuada para el almacén de este trabajo, pero en otros almacenes un operario puede entrar a un pasillo sin
salida y volver por donde vino, de modo que existen instancias físicamente
posibles que la formulación declara infactibles. Por ejemplo, si una ubicación requerida está al fondo de un pasillo sin salida que 
no contiene el depósito ni el punto de entrega, la única forma de visitarla es 
volver por el mismo camino, lo que obliga a pasar dos veces por la ubicación que 
da acceso al pasillo. Lo mismo puede ocurrir cuando un lote debe visitar casi todas las
ubicaciones de un almacén pequeño y no existe un camino simple que las recorra
a todas. Para cubrir esos casos habría que admitir rutas que repitan ubicaciones,
reescribiendo las restricciones de ruteo y el cálculo del tiempo de ruta. Es
posible, pero complica la formulación y la resolución, y no forma parte del
problema que se estudia.

{\bf Validez de $M^{\text{seq}}$ con espera voluntaria.} La cota de
$M^{\text{seq}}$ (Sección~\in[sec:notacion]) se basa en el tiempo total de
trabajo, lo que alcanza cuando cada lote termina en cuanto completa su
ejecución. Con espera voluntaria, en cambio, la separación entre las
finalizaciones de dos lotes de un mismo {\em picker} puede superar la cota, y
la restricción de secuencia que debería estar inactiva puede excluir la
solución óptima. El Anexo~\in[Anexo_Modelos] muestra un ejemplo y una cota
válida en todos los casos. Cuando la espera no conviene
(Sección~\in[sec:espera_voluntaria]), la cota original alcanza.

{\bf Lotes vacíos.} La formulación no distingue los lotes sin órdenes: si
$\mathcal{B}$ contiene alguno, igualmente le asigna una ruta del depósito al
nodo final, un {\em picker} y un lugar en una secuencia, y su finalización
entra en el objetivo. Además, los lotes vacíos son intercambiables entre sí,
lo que multiplica las soluciones equivalentes que debe considerar el
resolvedor. Como los lotes son un dato del problema y no una decisión de la
formulación, se considera que este caso debe resolverse antes, al conformar
los lotes.



\stopsection
\stopchapter

%========================================================================================
% GUÍA DE TRABAJO - CAPÍTULO 3 (Modelos lineales)
% Comentarios: no se compilan ni salen en el PDF. Borrar antes de la entrega.
% Armada el 21/10/2026 y revisada el 02/11/2026 sobre el texto actual. Se quitaron los
% ítems ya resueltos: los supuestos (congestión, capacidad, velocidad, cadena de frío,
% extremos de la ruta, camino simple y lotes vacíos), las condiciones de M y M_seq, el
% argumento de exactitud del Split, cómo se resuelve, el ciclo suelto por ubicaciones no
% requeridas, la reducción al modelo de referencia, el cuadro de diferencias, el nivel de
% servicio (Anexo C), la nota de las restricciones de demora y las dos decisiones
% (M_seq con espera quedó como característica de la formulación, con el Anexo C, y
% v_start distinto de v_end quedó como supuesto). Los TODO del texto siguen abiertos.
% Además: la descripción del modelo de referencia pasó a 2.6 (acá queda solo qué se agrega
% y qué se simplifica) y la M grande con varias binarias pasó a 2.1. Se decidió no
% incluir pseudocódigos de los modelos lineales: la Figura 3.1 cumple ese papel.
% [ ] = tarea; (D) = decisión de ustedes.
%========================================================================================
%
% 1. PENDIENTE

% (D) TODOs.
% [ ] Validación de los datos: se describe en 3.3.3 y la vuelven a listar 4.1.3, 4.2.3,
%     5.2 y 5.5. Al revisar esos capítulos, que remitan a 3.3.3.
% [ ] Paralelismo de la Etapa 1: queda en 3.3.1 (con su TODO) y en el cuadro de 3.5. El
%     detalle de cómo se ejecutó va en 5.4.
%
% 2. FUENTES
% - linear-picking-routing: models/opp_v5_1_mono/opp_model_v5_1_mono.zpl,
%   models/opp_v5_2_split/*.zpl, models/opp_v5_2_split/utils/cold_chain_validation.py,
%   tests/infra/zimpl/dat_files_builder.py (_compute_m_seq), tests/infra/zimpl/zpl_utils.py;
%   docs/thesis/models.tex (secs. 2 a 6).
% - Notas/opp_v5_development.tex: secs. "El problema y sus supuestos" a "Relación entre las
%   dos formulaciones", "Implementación" y "Estudio 07".
%========================================================================================
% NOTA: Añadí todo, pero con comentarios para revisar. No creo que esté finalizado,
% creo que hay cosas que conviene pulir antes de cerrar este capítulo.
